\begin{itemize}
   \item When $ M > N $, $\log_a M> \log_aN$; for example $100 > 10$, but $\log_{10} 100 = 2, \log_{10} 10 =1$ so $\log_{10} 100 > \log_{10} 10$.
   \item $\log_a MN = \log_a M +\log_a N$; for example $\log_2 8 = 3,\log_2 16 = 4$, but by the definition of logarithms this represents the relationship $2^3=8,2^4=16$. Here $\log_2 8 \times 16$ is the $\log_2$ of $2^3 \times 2^4$, which equals $2^{(3+4)}=2^7$, so we have $\log_2 8 \times 16 = \log_2 8 + \log_2 16$, where multiplication becomes addition.
   \item $\log_a \frac{M}{N} = \log_a M - \log_a N$; similarly, $\log_2 \frac{16}{8}=\log_2 2=1$, and we can compute $\log_2 16 - \log_2 8 = 3-2=1$.
\end{itemize}
